\chapter*{Phụ Lục VII. Lời Gửi Các Em Học Sinh}
\addcontentsline{toc}{chapter}{Phụ Lục VII. Lời Gửi Các Em Học Sinh}
\markboth{Lời Gửi Các Em Học Sinh}{Lời Gửi Các Em Học Sinh}

\begin{truyen}{Thầy Không Muốn Các Em Học Quá Muộn}{I Do Not Want You to Learn Too Late}
\chuhoa{T}{hầy viết thêm những trang này cho học sinh, bởi có những bài học nếu để đến khi làm cha, làm mẹ, làm người gánh cả gia đình rồi mới hiểu thì thường rất đau.}
\textit{(I write these pages for students because some lessons become very painful if they are only understood after one becomes a parent, a provider, a person carrying an entire family.)}

Khi còn đi học, các em dễ nghĩ chuyện tiền bạc là việc của người lớn. Nhưng chính lúc còn trẻ, con người đang tập hình thành thói quen: tiêu theo cảm xúc hay biết dừng lại, so sánh với bạn bè hay biết sống theo sức mình, dùng đồ có biết quý hay coi mọi thứ là hiển nhiên. Những thói quen ấy lớn lên cùng con người nhanh hơn chính các em tưởng.
\textit{(While in school, you may think money is an adult matter. Yet youth is exactly the time when habits are formed: whether to spend emotionally or pause, whether to compare yourself with friends or live within your means, whether to value what you use or treat everything as automatic. Such habits grow with a person faster than students often imagine.)}

Thầy không mong các em trở nên tính toán hẹp hòi. Thầy chỉ mong các em hiểu rằng sống biết quý tiền là một phần của đạo đức. Ai biết quý công sức, quý bữa cơm, quý sách vở, quý thời gian của cha mẹ và quý cả lao động của mình, người ấy sau này sẽ khó phung phí cuộc đời hơn.
\textit{(I do not hope you become narrow or miserly. I only hope you understand that valuing money is part of moral character. Whoever values effort, meals, books, the time of parents, and one’s own labor will later find it harder to waste life itself.)}
\end{truyen}

\begin{baihoc}
Dạy học sinh về tiền không phải để các em ám ảnh với tiền. Đó là để các em biết ơn hơn, tự trọng hơn và trưởng thành sớm hơn trong cách sống.
\textit{(Teaching students about money is not to make them obsessed with money. It is to help them become more grateful, more self-respecting, and more mature in how they live.)}
\end{baihoc}

\section*{Mười Hai Điều Thầy Muốn Học Sinh Ghi Nhớ}
\begin{truyen}{Những Điều Nhỏ Nhưng Đi Rất Xa}{Small Lessons That Travel Far}
1. Đừng tập tiêu tiền để được công nhận.\
\textit{(Do not train yourself to spend in order to be recognized.)}

2. Một cây bút, một cuốn tập, một bữa cơm đều có công sức ở sau.\
\textit{(A pen, a notebook, and a meal all carry labor behind them.)}

3. Tiết kiệm không đối nghịch với lòng tốt.\
\textit{(Saving is not the opposite of kindness.)}

4. Biết quý đồ dùng là dấu hiệu của người biết quý cuộc sống.\
\textit{(Valuing your belongings is a sign of valuing life.)}

5. So sánh với bạn bè là một con đường ngắn dẫn đến bất mãn.\
\textit{(Comparing yourself to friends is a short road to discontent.)}

6. Cha mẹ không chỉ cho tiền; cha mẹ đang cho cả thời gian đời mình.\
\textit{(Parents do not only give money; they give pieces of their life.)}

7. Càng tiêu để khoe, càng dễ nghèo bên trong.\
\textit{(The more you spend to show off, the easier it is to become inwardly poor.)}

8. Một thói quen tốt nhỏ đáng giá hơn một lời hứa lớn.\
\textit{(One small good habit is worth more than one grand promise.)}

9. Đừng để nỗi buồn điều khiển ví tiền.\
\textit{(Do not let sadness control your wallet.)}

10. Muốn trưởng thành, phải biết chậm lại trước cám dỗ.\
\textit{(If you want to mature, learn to slow down before temptation.)}

11. Người biết đủ thường học và sống sâu hơn.\
\textit{(Those who know enough often study and live more deeply.)}

12. Giữ tiền đúng cũng là giữ lòng tự trọng của mình.\
\textit{(Keeping money wisely is also keeping your self-respect.)}
\end{truyen}

\section*{Mười Hai Câu Hỏi Thầy Có Thể Dùng Trên Lớp}
\begin{truyen}{Để Một Tiết Học Không Chỉ Dừng Ở Kiến Thức}{So That a Class Does Not End with Information Alone}
1. Em có đang muốn một món đồ vì thật sự cần hay vì thấy bạn mình có?\
\textit{(Do you want an item because you need it or because your friend has it?)}

2. Em đã bao giờ làm cha mẹ buồn vì tiêu xài vô tâm chưa?\
\textit{(Have you ever made your parents sad through careless spending?)}

3. Em quý nhất điều gì cha mẹ đang hy sinh cho mình?\
\textit{(What sacrifice from your parents do you value most?)}

4. Điều gì trong lớp học cho thấy một người có nếp sống chừng mực?\
\textit{(What in the classroom reveals a person of moderation?)}

5. Tại sao biết đủ lại liên quan đến hạnh phúc?\
\textit{(Why is knowing enough related to happiness?)}

6. Vì sao người tiêu tiền để khoe thường khó bình yên lâu dài?\
\textit{(Why do people who spend to impress often struggle to remain at peace?)}

7. Tiết kiệm có phải là keo kiệt không? Vì sao?\
\textit{(Is saving the same as stinginess? Why or why not?)}

8. Một học sinh biết quý sách vở sẽ khác gì với một học sinh coi thường đồ dùng?\
\textit{(How is a student who values books different from one who takes belongings lightly?)}

9. Em đang lãng phí điều gì ngoài tiền?\
\textit{(What are you wasting besides money?)}

10. Nếu phải tự răn mình bằng một câu, em sẽ viết câu gì?\
\textit{(If you had to discipline yourself with one sentence, what would it be?)}

11. Theo em, đạo đức có liên quan đến cách dùng tiền không?\
\textit{(Do you think morality is related to the way one uses money?)}

12. Bài học nào em muốn mang về nhà sau khi đọc quyển sách này?\
\textit{(What lesson do you want to take home after reading this book?)}
\end{truyen}

\begin{ghinhoanh}
Students do not need a lecture first.

They need a truthful example and a clear question.
\end{ghinhoanh}

\ngancach